%=====================================================================
%  SUPPLEMENTARY MATERIAL
%  Threat-Model Dependence in the Adversarial Robustness of
%  Mamba-Based Network Intrusion Detection
%
%  This document holds the material the main paper refers to as
%  released with the archived artefact: the derivations and
%  pseudo-code omitted from Section 3, the extended scope statement
%  omitted from Section 4, and the extended leakage, padding-mask,
%  severity, defence-verification, threat-model and capacity analyses
%  omitted from Section 5.
%
%  Every quantity below is reproduced from the released result files;
%  no value is transcribed by hand.  The provenance of each is named
%  in the text.
%=====================================================================
\documentclass[11pt,a4paper]{article}
\usepackage[utf8]{inputenc}
\usepackage[T1]{fontenc}
\usepackage[english]{babel}
\usepackage{mathptmx}
\usepackage[a4paper,margin=2.5cm]{geometry}
\usepackage{amsmath,amssymb,amsfonts,bm}
\usepackage{booktabs}
\usepackage{array}
\usepackage{multirow}
\usepackage[font=small,labelsep=period,labelfont=bf]{caption}
\usepackage[ruled,vlined,linesnumbered]{algorithm2e}
\usepackage{url}
\usepackage[hyphens]{xurl}
\sloppy
\emergencystretch=2em
\usepackage[hidelinks,breaklinks]{hyperref}
\SetAlgorithmName{Algorithm}{Algorithm}{List of Algorithms}
\SetKwInput{KwData}{Input}
\SetKwInput{KwResult}{Output}

\newcommand{\R}{\mathbb{R}}
\newcommand{\E}{\mathbb{E}}
\newcommand{\Ind}{\mathbb{1}}
\newcommand{\vx}{\bm{x}}
\newcommand{\vy}{\bm{y}}
\newcommand{\vh}{\bm{h}}
\newcommand{\vu}{\bm{u}}
\newcommand{\vz}{\bm{z}}
\newcommand{\vdelta}{\bm{\delta}}
\newcommand{\vtheta}{\bm{\theta}}
\newcommand{\mA}{\bm{A}}
\newcommand{\mB}{\bm{B}}
\newcommand{\mC}{\bm{C}}
\newcommand{\mD}{\bm{D}}
\newcommand{\Abar}{\overline{\bm{A}}}
\newcommand{\Bbar}{\overline{\bm{B}}}
\newcommand{\Dset}{\mathcal{D}}
\newcommand{\Xset}{\mathcal{X}}
\newcommand{\Sset}{\mathcal{S}}
\DeclareMathOperator*{\argmin}{arg\,min}
\DeclareMathOperator*{\argmax}{arg\,max}
\DeclareMathOperator{\softmax}{softmax}
\DeclareMathOperator{\softplus}{softplus}
\DeclareMathOperator{\SiLU}{SiLU}
\DeclareMathOperator{\LN}{LayerNorm}
\DeclareMathOperator{\sign}{sign}
\DeclareMathOperator{\diag}{diag}
\DeclareMathOperator{\clip}{clip}
\DeclareMathOperator{\Proj}{\Pi}

\renewcommand{\thesection}{S\arabic{section}}
\renewcommand{\theequation}{S\arabic{equation}}
\renewcommand{\thetable}{S\arabic{table}}

\title{Supplementary Material\\[2pt]
\large Threat-Model Dependence in the Adversarial Robustness of
Mamba-Based Network Intrusion Detection}
\author{Haydar Kilic\\[2pt]
\normalsize Department of Artificial Intelligence Engineering,
Ostim Technical University, Ankara, Turkey\\
\normalsize \texttt{haydar.kilic@ostimteknik.edu.tr}}
\date{}

\begin{document}
\maketitle
\vspace{-2.5em}

\noindent
This document holds the material the main paper refers to as released
with the archived artefact. Section~\ref{sup:deriv} gives the
derivations and definitions omitted from the paper's methodology
section and Section~\ref{sup:algo} the two pseudo-code listings;
Section~\ref{sup:scope} extends the statement of scope;
Sections~\ref{sup:leak}--\ref{sup:capacity} give in full the analyses
the paper reports only in summary. Every quantity is reproduced from
the result files deposited with the artefact, and the file each comes
from is named where it is used. Section~\ref{sup:index} indexes the
deposited tables and figures.

\medskip\noindent
The main paper cites the sections of this document as S1--S10. Code,
result frames and this document are archived together at
\url{https://doi.org/10.5281/zenodo.23497561} (concept identifier;
release v2.0.1); the corpus subsamples are in the companion record
\url{https://doi.org/10.5281/zenodo.23499680}. This document is not
self-contained and is to be read with the paper.

\tableofcontents

%=====================================================================
\section{Derivations and definitions omitted from the paper}
\label{sup:deriv}

\subsection{Class-weighted empirical risk}

Operational NIDS traffic is dominated by benign records, and an
unweighted empirical risk admits a minimiser that predicts the majority
class unconditionally. The detector parameters are therefore estimated
by minimising the class-weighted risk
\begin{equation}
\vtheta^{\star}
=\argmin_{\vtheta}\;
\frac{1}{|\Dset|}\sum_{i=1}^{|\Dset|} w_{y^{(i)}}\,
\ell\bigl(f_{\vtheta}(\vx^{(i)}),\,y^{(i)}\bigr),
\label{eq:erm}
\end{equation}
with $\ell$ the categorical cross-entropy and the weights the
inverse-frequency coefficients
\begin{equation}
w_{k}=\frac{|\Dset|}{K\,n_{k}},\qquad
n_{k}=\sum_{i=1}^{|\Dset|}\Ind\bigl[y^{(i)}=k\bigr],\quad K=2 .
\label{eq:weights}
\end{equation}
Equation~(\ref{eq:weights}) is the exact analogue of the balanced
class-weight convention adopted for the Random Forest baseline, so the
imbalance correction is identical across the deep and non-deep arms of
the comparison. The predictive distribution and decision rule are
\begin{equation}
p_{\vtheta}(y\mid\vx)=\softmax\bigl(f_{\vtheta}(\vx)\bigr),
\qquad
\hat{y}(\vx)=\argmax_{k\in\{0,1\}} \; f_{\vtheta}(\vx)_{k}.
\label{eq:decision}
\end{equation}

\subsection{Standardisation and the projection operator}

Continuous features are standardised with statistics estimated
exclusively on the training partition and applied unchanged to the test
partition,
\begin{equation}
\tilde{x}_{j}=\frac{x_{j}-\mu_{j}^{\mathrm{tr}}}
{\sigma_{j}^{\mathrm{tr}}+\epsilon},\qquad \epsilon=10^{-8},
\label{eq:standardise}
\end{equation}
and the Euclidean projection onto the feasible set of the paper's
Equation~(1) is applied component-wise,
\begin{equation}
\Proj_{\Sset_{\varepsilon}(\vx)}(\vz)=
\clip\Bigl(\vx+\clip(\vz-\vx,\,-\varepsilon,\,\varepsilon),\;l,\;u\Bigr),
\label{eq:proj}
\end{equation}
where $l=\min_{\Dset_{\mathrm{tr}}}\vx$ and
$u=\max_{\Dset_{\mathrm{tr}}}\vx$ are the empirical bounds of the
training distribution in standardised feature space.

\subsection{Continuous-time formulation and zero-order-hold discretisation}

A linear time-invariant state-space model maps a scalar input $u(t)$ to
an output $y(t)$ through a latent state $\vh(t)\in\R^{N}$:
\begin{equation}
\frac{\mathrm{d}\vh(t)}{\mathrm{d}t}=\mA\vh(t)+\mB u(t),
\qquad
y(t)=\mC\vh(t)+D\,u(t),
\label{eq:ssm-cont}
\end{equation}
with $\mA\in\R^{N\times N}$, $\mB\in\R^{N\times1}$,
$\mC\in\R^{1\times N}$ and a skip coefficient $D\in\R$. Discretising
(\ref{eq:ssm-cont}) with a zero-order hold over a step size $\Delta>0$
yields
\begin{align}
\vh_{t}&=\Abar\vh_{t-1}+\Bbar u_{t}, \label{eq:ssm-disc-h}\\
y_{t}&=\mC\vh_{t}+D\,u_{t}, \label{eq:ssm-disc-y}
\end{align}
where, following the reference formulation of Gu and Dao with a
diagonal $\mA$,
\begin{equation}
\Abar=\exp(\Delta\mA),
\qquad
\Bbar=\Delta\mB.
\label{eq:zoh}
\end{equation}
The expression for $\Bbar$ is the first-order term of the exact
zero-order-hold coefficient
$(\Delta\mA)^{-1}\bigl(\exp(\Delta\mA)-\bm{I}\bigr)\Delta\mB$, whose
expansion is $\Delta\mB+O(\Delta^{2})$; the truncation is adopted for
numerical stability and for consistency with the reference
implementation. The state matrix is parameterised as
\begin{equation}
\mA=-\exp(\mA_{\log}),\qquad
\mA_{\log}\in\R^{d_{\mathrm{inner}}\times N},
\label{eq:Alog}
\end{equation}
so that every entry of $\mA$ is strictly negative and, since
$\Delta>0$, every entry of $\Abar$ lies in $(0,1)$: the recurrence is
contractive, which precludes the exploding-state pathology of naive
recurrent parameterisations. Under the S4D-real initialisation the
$n$-th diagonal entry is set to $\mA_{\log}=\log n$ for
$n=1,\dots,N$, giving $\mA\in\{-1,\dots,-N\}$; this is the sequence the
source writes as $A_{n}=-(n+1)$ under zero-based indexing. With $N=16$
the resulting decay rates span a sixteenfold range, so the state
channels integrate the input over time-scales differing by more than an
order of magnitude.

\subsection{The selective projection and the block dimensions}

The selective parameters are obtained by a single projection of the
internal representation $\vu_{t}\in\R^{d_{\mathrm{inner}}}$:
\begin{equation}
\bigl[\tilde{\vdelta}_{t},\,\mB_{t},\,\mC_{t}\bigr]
=\bm{W}_{x}\vu_{t},
\label{eq:selproj}
\end{equation}
whose three blocks have widths $d_{\mathrm{inner}}$, $N$ and $N$
respectively, so that
$\tilde{\vdelta}_{t}\in\R^{d_{\mathrm{inner}}}$,
$\mB_{t},\mC_{t}\in\R^{N}$ and
$\bm{W}_{x}\in\R^{(d_{\mathrm{inner}}+2N)\times d_{\mathrm{inner}}}$,
with
$\bm{W}_{\Delta}\in\R^{d_{\mathrm{inner}}\times d_{\mathrm{inner}}}$.
In the block equations of the paper the affine maps
$\bm{W}_{\mathrm{in}}\in\R^{2d_{\mathrm{inner}}\times d_{\mathrm{model}}}$
and $\bm{W}_{\mathrm{out}}$ are applied row-wise, that is,
independently at each step, whereas the depth-wise causal convolution
and the selective scan act along the sequence axis; the gating endows
the block with a multiplicative input-dependent pathway and the
residual is post-normalised, the normalisation being applied after the
residual is added.

\subsection{Detector head}

The detector embeds the univariate pseudo-sequence, adds a learned
positional embedding $\bm{P}\in\R^{L\times d_{\mathrm{model}}}$
initialised from a truncated normal distribution with $\sigma=0.02$,
applies $K=3$ stacked blocks, mean-pools over the sequence axis and
classifies:
\begin{align}
\vh^{(0)}&=\vx\bm{W}_{e}+\bm{P},\qquad
\bm{W}_{e}\in\R^{1\times d_{\mathrm{model}}},\\
\vh^{(k)}&=\mathrm{S6Block}_{k}\bigl(\vh^{(k-1)}\bigr),
\quad k=1,\dots,K,\\
\bar{\vh}&=\frac{1}{L}\sum_{t=1}^{L}\vh^{(K)}_{t},\\
f_{\vtheta}(\vx)&=\bm{W}_{2}\,\mathrm{GELU}
\bigl(\bm{W}_{1}\LN(\bar{\vh})\bigr),
\end{align}
with
$\bm{W}_{1}\in\R^{(d_{\mathrm{model}}/2)\times d_{\mathrm{model}}}$
and $\bm{W}_{2}\in\R^{2\times(d_{\mathrm{model}}/2)}$.

\subsection{Closed-form evaluation used in the second validation check}

For diagonal $\mA$ the time-varying recurrence admits the closed form
\begin{equation}
\vh_{t}=\sum_{s=1}^{t}
\Bigl(\textstyle\prod_{r=s+1}^{t}\Abar_{r}\Bigr)
\odot\Bbar_{s}\otimes\vu_{s},
\label{eq:closedform}
\end{equation}
in which the empty product at $s=t$ is the identity. Because
$\Abar_{r}=\exp(\Delta_{r}\otimes\mA)$, the product
in~(\ref{eq:closedform}) is evaluated without a recurrence as
$\exp(\bm{S}_{t}-\bm{S}_{s})$, where
$\bm{S}_{t}=\sum_{r\le t}\Delta_{r}\otimes\mA$ is a cumulative sum in
log space; since $\mA\le 0$ and $\Delta_{r}>0$, every such factor lies
in $(0,1]$ and the evaluation is numerically stable. This is the
reference against which the sequential loop is compared in the paper's
second validation check, which reports a maximum deviation of
$1.39\times10^{-17}$ in double precision
(\texttt{results/verification/verification.json}, field
\texttt{check2}~$=1.3877787807814457\times10^{-17}$).

\subsection{Optimisation schedule}

The learning rate is annealed by a cosine schedule stepped once per
epoch,
\begin{equation}
\eta_{e}=\tfrac{1}{2}\eta_{0}
\Bigl(1+\cos\bigl(\pi e/N_{\mathrm{ep}}\bigr)\Bigr),
\qquad N_{\mathrm{ep}}=40,
\label{eq:cosine}
\end{equation}
the epoch budget being written $N_{\mathrm{ep}}$ to avoid collision
with the expansion factor $E$.

Model selection and early stopping are governed by the balanced
accuracy on the validation split, which the paper states in words and
which is
\begin{equation}
\mathrm{BA}=\frac{1}{K}\sum_{k=1}^{K}
\frac{\sum_{i}\Ind[\hat{y}^{(i)}=k]\,\Ind[y^{(i)}=k]}
{\sum_{i}\Ind[y^{(i)}=k]},
\qquad K=2,
\label{eq:balacc}
\end{equation}
the unweighted mean of the per-class recalls. It is used in place of
raw accuracy because a degenerate majority-class predictor attains a
raw accuracy equal to the majority prior while attaining
$\mathrm{BA}=0.5$, and because it is the quantity the class weights of
Equation~(\ref{eq:weights}) make the training objective an analogue of,
so selection and optimisation target the same criterion.

\subsection{Per-category detection rates}

For every attack category $k$ with at least five test instances the
clean and adversarial detection rates are
\begin{align}
D_{k}^{\mathrm{clean}}&=\frac{1}{|I_{k}|}\sum_{i\in I_{k}}
\Ind\bigl[\hat{y}(\vx_{i})=1\bigr],
\label{eq:detrate-clean}\\
D_{k}^{\mathrm{adv}}&=\frac{1}{|I_{k}|}\sum_{i\in I_{k}}
\Ind\bigl[\hat{y}(\vx^{\mathrm{adv}}_{i})=1\bigr],
\label{eq:detrate-adv}
\end{align}
where $I_{k}=\{i: y_{i}=1,\,c_{i}=k\}$, and the evasion success rate of
the paper's Equation~(16) is their difference. The analogous
quantities for the two binary classes are computed by the same
procedure and are released with the artefact rather than tabulated,
since the category-level decomposition subsumes them for the malicious
class.

\subsection{Latency and throughput}

Per-sample latency and throughput follow from the measured per-batch
wall-clock time $\bar{\lambda}_{\mathrm{batch}}$ in milliseconds at
batch size $b$ as
\begin{equation}
\lambda_{\mathrm{sample}}=\frac{\bar{\lambda}_{\mathrm{batch}}}{b},
\qquad
\mathrm{TP}=\frac{1000\,b}{\bar{\lambda}_{\mathrm{batch}}} .
\label{eq:latency}
\end{equation}
Measurement uses CUDA events with explicit device synchronisation after
$20$ warm-up iterations over $100$ timed repetitions. The Random Forest
baseline executes on CPU and is timed with a host-side clock over a
reduced schedule of five warm-up and twenty timed iterations, so its
figures are reported for reference only.

\subsection{Paired testing and the Holm--Bonferroni correction}

Because the seed is shared between the proposed model and its
comparator, the comparison is paired, and the parametric statistic is
\begin{equation}
t=\frac{\bar{d}}{s_{d}/\sqrt{S}},
\qquad
d_{s}=a_{s}^{\mathrm{Mamba}}-a_{s}^{\mathrm{comp}},
\label{eq:ttest}
\end{equation}
over $S=7$ seeds. Since one hypothesis is tested per (dataset, attack)
cell, the family-wise error rate over the resulting family of twelve is
controlled by the Holm--Bonferroni step-down procedure: with ordered
raw $p$-values $p_{(1)}\le\dots\le p_{(m)}$, the adjusted values are
\begin{equation}
\tilde{p}_{(j)}=\max_{i\le j}\;
\min\bigl\{1,\;(m-i+1)\,p_{(i)}\bigr\},
\label{eq:holm}
\end{equation}
and $H_{(j)}$ is rejected iff $\tilde{p}_{(j)}\le\alpha=0.05$. The two
conventions the paper describes differ only in the choice of
comparator; the adjusted values therefore differ between them even
where the raw $p$-value is identical, because the correction is applied
within each family separately. The clearest instance is NF-UNSW-NB15
under transfer, where the raw $p$-value of $0.0223$ becomes
$\tilde{p}=0.1558$ in the pre-specified family and $\tilde{p}=0.1113$
in the post-hoc family
(\texttt{results/tables/T2\_istatistiksel\_anlamlilik.csv} and
\texttt{T2b\_istatistiksel\_anlamlilik\_posthoc.csv}). Every
significance claim in the paper is made with respect to the
pre-specified family.

%=====================================================================
\section{Pseudo-code}
\label{sup:algo}

\begin{algorithm}[H]
\small
\caption{Leakage-aware partitioning. The strongest guarantee supported
by the metadata present in a given corpus is selected: a group split
when flow or endpoint identifiers exist, a chronological split when a
parsable timestamp exists, and a stratified random split only as a last
resort, in which case a leakage caveat is written to the run log.}
\label{alg:split}
\KwData{Records $\Dset$; candidate group columns $G$; timestamp column
$\tau$; test fraction $\rho_{\mathrm{te}}=0.2$; fixed split seed
$s=42$, shared by all experimental seeds}
\KwResult{Partition assignment $\pi\in\{\text{train},\text{test}\}^{N}$}
\BlankLine
$G^{\ast}\leftarrow$ columns of $G$ present in $\Dset$\;
\uIf{$G^{\ast}\neq\emptyset$}{
  $k\leftarrow$ composite key from the first two columns of $G^{\ast}$\;
  \If{$|\mathrm{unique}(k)|\ge 10$}{
    $\pi\leftarrow\textsc{GroupShuffleSplit}(\Dset,k,\rho_{\mathrm{te}},s)$\;
    \Return $\pi$ \tcp*{no group spans both partitions}
  }
}
\uIf{$\tau$ present \textbf{and} $\ge 90\%$ of $\tau$ parses}{
  sort indices by $\tau$ (stable)\;
  assign the latest $\lceil\rho_{\mathrm{te}}N\rceil$ records to \textit{test}\;
  \Return $\pi$ \tcp*{chronological hold-out}
}
$\pi\leftarrow\textsc{StratifiedSplit}(\Dset,y,\rho_{\mathrm{te}},s)$\;
record an explicit leakage caveat in the run log\;
\Return $\pi$\;
\end{algorithm}

\begin{algorithm}[H]
\small
\caption{Selective scan (S6) forward pass for one block. $B$ is the
batch size, $L$ the sequence length, $N$ the state dimension and
$d_{\mathrm{inner}}=E\,d_{\mathrm{model}}$ the expanded channel width.
$\odot$ is the Hadamard product and $\otimes$ broadcasts along the
state dimension. Lines~7--10 realise the time-varying recurrence as an
explicit loop.}
\label{alg:s6}
\KwData{$\vx\in\R^{B\times L\times d_{\mathrm{model}}}$;
parameters $\bm{W}_{\mathrm{in}},\bm{W}_{x},\bm{W}_{\Delta},
\bm{W}_{\mathrm{out}},\mA_{\log},\mD$, the depth-wise convolution
kernel $\bm{K}\in\R^{d_{\mathrm{inner}}\times k}$ and the
$\LN$ scale and shift}
\KwResult{$\mathrm{out}\in\R^{B\times L\times d_{\mathrm{model}}}$}
\BlankLine
$[\vu',\vz]\leftarrow\bm{W}_{\mathrm{in}}\vx$\;
$\vu\leftarrow\SiLU\bigl(\mathrm{DWConv}_{k=4}(\vu')[:,{:}L]\bigr)$
\tcp*{causal}
$[\tilde{\vdelta},\mB,\mC]\leftarrow\bm{W}_{x}\vu$\;
$\Delta\leftarrow\softplus(\bm{W}_{\Delta}\tilde{\vdelta})$;\quad
$\mA\leftarrow-\exp(\mA_{\log})$\;
$\Abar\leftarrow\exp(\Delta\otimes\mA)$;\quad
$\Bbar\leftarrow\Delta\otimes\mB$\;
$\vh\leftarrow\bm{0}\in\R^{B\times d_{\mathrm{inner}}\times N}$\;
\For{$t\leftarrow1$ \KwTo $L$}{
  $\vh\leftarrow\Abar_{:,t}\odot\vh+\Bbar_{:,t}\otimes\vu_{:,t}$\;
  $\vy_{t}\leftarrow\sum_{n=1}^{N}\vh_{:,:,n}\,\mC_{:,t,n}$\;
}
$\vy\leftarrow\mathrm{stack}(\vy_{1},\dots,\vy_{L})+\mD\odot\vu$\;
$\vy\leftarrow\vy\odot\SiLU(\vz)$\;
\Return $\LN(\bm{W}_{\mathrm{out}}\vy+\vx)$\;
\end{algorithm}

%=====================================================================
\section{Extended statement of scope}
\label{sup:scope}

The complete experimental grid comprises three datasets, five
architectures, seven random seeds and ten attack settings (FGSM and PGD
at $\varepsilon\in\{0.01,0.03,0.05,0.1\}$ and black-box transfer at
$\varepsilon\in\{0.03,0.1\}$). Because the Random Forest is excluded
from all ten by construction, the master frame holds $945$ evaluation
records rather than the $1{,}155$ of a full product: $105$ clean cells
over five architectures and $840$ attacked cells over the four
differentiable ones.

The official NSL-KDD files were re-counted directly to confirm the
protocol statement the paper makes about them: \texttt{KDDTrain+} holds
$125{,}973$ records carrying $23$ distinct labels and
\texttt{KDDTest+} holds $22{,}544$ records carrying $38$, of which $17$
labels occur in the test file and in no training record. Those $17$ are
\texttt{apache2}, \texttt{httptunnel}, \texttt{mailbomb},
\texttt{mscan}, \texttt{named}, \texttt{processtable}, \texttt{ps},
\texttt{saint}, \texttt{sendmail}, \texttt{snmpgetattack},
\texttt{snmpguess}, \texttt{sqlattack}, \texttt{udpstorm},
\texttt{worm}, \texttt{xlock}, \texttt{xsnoop} and \texttt{xterm},
and they account for $3{,}750$ test records, $16.63\%$ of the
partition. The difference $38-23$ does not give $17$ because two
training labels (\texttt{spy} and \texttt{warezclient}) are absent from
the test file; the count of unseen types is therefore a set difference
and not a subtraction of cardinalities.

After sanitisation and tensorisation the three corpora comprise
$125{,}973$/$22{,}544$ (NSL-KDD, official split),
$111{,}975$/$27{,}289$ (NF-UNSW-NB15) and $250{,}143$/$62{,}370$
(CSE-CIC-IDS2018) training/test records, with malicious proportions of
$0.465$/$0.569$, $0.503$/$0.490$ and $0.529$/$0.500$. The record counts
written by the split verification
(\texttt{results/verification/verification.json}) precede the
sanitisation stage, which removes the non-finite rows: for
CSE-CIC-IDS2018 that file reports $250{,}918$/$62{,}730$, which
sanitisation reduces to the $250{,}143$/$62{,}370$ above.

The native feature dimensionalities are $41$, $41$ and $79$, so the
fixed pseudo-sequence length $L=64$ entails zero-padding of the
trailing $23$ coordinates, $35.94\%$ of the sequence, for the first two
corpora and the truncation of $15$ trailing attributes for
CSE-CIC-IDS2018. The automated health audit reported no residual
NaN/Inf entries in any partition and identified $24$, $23$ and $8$
near-zero-variance features respectively. These counts are largely an
artefact of the pseudo-sequence construction rather than of the
underlying telemetry: for NF-UNSW-NB15 all $23$ are padded positions,
and for NSL-KDD $23$ of the $24$ are, the remaining one being a single
genuinely constant input attribute. CSE-CIC-IDS2018 is truncated rather
than padded and therefore contributes no padded positions at all, so
all $8$ of the columns flagged on that corpus are attributes that are
genuinely constant across its training partition.

Counting a category as analysable when at least five instances survive
in the test partition retained four categories for NSL-KDD (DoS, Probe,
R2L, U2R), nine for NF-UNSW-NB15 and two for CSE-CIC-IDS2018 (Bot and
Infiltration, the latter spelled \texttt{Infilteration} in the released
label column). The last of these three counts is a consequence of the
partitioning rather than of the corpus: the file carries fourteen
distinct attack labels, but the chronological cut distributes them
almost disjointly, twelve occurring only within the training interval,
Infiltration alone being present on both sides of the cut, and Bot
occurring exclusively in the test partition. Detection on that corpus
is thus partly a zero-shot task.

Neither the group-split check nor the chronological-split check can be
reproduced from the released tensors, since the address and timestamp
columns are dropped before tensorisation; both are therefore performed
from the raw corpora and their output is released in
\texttt{results/verification/}. The group split on
\texttt{IPV4\_SRC\_ADDR} and \texttt{IPV4\_DST\_ADDR} resolved to $242$
unique host pairs with zero groups shared between the two partitions,
and the timestamp column of CSE-CIC-IDS2018 parsed for every record
(\texttt{parsed\_fraction}~$=1.0$) with the training and test intervals
confirmed disjoint.

%=====================================================================
\section{The NF-UNSW-NB15 leakage episode in full}
\label{sup:leak}

\subsection{The defect and the route by which it was found}

The released NF-UNSW-NB15-v2 CSV carries $45$ columns, of which two are
labels: the binary \texttt{Label} and \texttt{Attack}, the attack-family
name. The preprocessing stage of the first execution excluded the
former but not the latter, and because \texttt{Attack} is a text column
it was label-encoded and admitted as a feature. Its value set for benign
flows is disjoint from its value set for malicious ones, so a
membership rule on that one coordinate reproduces the label exactly, at
an accuracy of $1.0000$ on both partitions.

The single-threshold audit did not catch it, and the reason is exact
rather than incidental. Alphabetical encoding of the ten values places
\texttt{Benign} third, as code $2$ behind \texttt{Analysis} and
\texttt{Backdoor}, so every single cut assigns at least one class
wrongly: the best of them, $x>2$, misclassifies the $2{,}299$
\texttt{Analysis} and $2{,}169$ \texttt{Backdoor} flows as benign and
therefore scores $1-4{,}468/139{,}264=0.9679$ on the released file,
against the $1.0000$ a membership test returns. Both figures are
reproducible from the deposited subsample:

\begin{verbatim}
d = pd.read_csv("NF-UNSW-NB15-v2.csv")
cats = sorted(d.Attack.unique())
x = d.Attack.map({c: i for i, c in enumerate(cats)})
((x != cats.index("Benign")) == d.Label).mean()   # 1.0
((x > 2) == d.Label).mean()                       # 0.967917
\end{verbatim}

\noindent A gap of $0.032$ between the two audits is what separated a
corpus that looked merely saturated from one in which a single
coordinate carried the label. The defect
was found instead by comparing the feature count the pipeline reported
against the count the feature-importance analysis used, and the corpus
was reprocessed and retrained with \texttt{Attack} withheld. Every
NF-UNSW-NB15 figure in the paper comes from that corrected execution;
NSL-KDD and CSE-CIC-IDS2018 were never affected and were verified
record by record to be identical to their original values.

The episode is reported rather than silently repaired because the
detection route matters. A membership test over per-class value sets
catches this failure and the threshold test that is conventional in
leakage audits does not, so an audit built only on the latter will pass
a corpus in which one coordinate determines the label.

\subsection{Why saturation persists after the correction}

That a $200$-tree Random Forest attains $0.9976$ accuracy with a
seed dispersion of $2.5\times10^{-5}$ is the pattern one expects from a
dominant, near-deterministic feature rather than from a hard problem
solved well. We therefore examined the impurity-based feature
importances of that model together with the best accuracy obtainable
from each feature under a single threshold
(\texttt{results/verification/verification.json}, field
\texttt{leakage}). Three features are individually almost sufficient:
\texttt{MAX\_TTL} and \texttt{MIN\_TTL} each separate the two classes
at $0.9948$ accuracy with a single threshold and
\texttt{MIN\_IP\_PKT\_LEN} at $0.9916$, and together they carry $0.576$
of the total impurity-based importance mass.

The time-to-live fields are the more diagnostic of the two kinds. A TTL
value is a function of the operating system that emitted the packet and
of the number of hops it traversed, not of whether the traffic is
hostile; that it predicts the label at $99.5\%$ indicates that benign
and malicious traffic in this capture originate from hosts that are
distinguishable by network position or platform. Predictive power of
this kind is a property of the testbed rather than of intrusion, and
partitioning cannot remove it, since it is carried in the feature
vector itself.

\subsection{The feature-removal ablation}

Whether the models rely on these features is a separate question, and
it was tested directly rather than inferred. Table~\ref{tab:ablation}
reports the result, computed on the group-aware split
(\texttt{results/verification/ttl\_ablation.json}).

\begin{table}[h]
\centering
\caption{Feature-removal ablation on NF-UNSW-NB15. Held-out accuracy of
the Random Forest under the group-aware split, with the
near-deterministic features withheld in turn. The $41$ features of the
all-features row are the $43$ non-label columns of the released CSV
less the two endpoint-address fields reserved as group keys.}
\label{tab:ablation}
\small
\begin{tabular}{@{}lccc@{}}
\toprule
Features withheld & $d$ & Accuracy & Change \\
\midrule
none                                   & 41 & 0.997618 & --- \\
\texttt{MAX\_TTL}                      & 40 & 0.997655 & $+3.7\times10^{-5}$ \\
\texttt{MAX\_TTL}, \texttt{MIN\_TTL}   & 39 & 0.997508 & $-1.1\times10^{-4}$ \\
\;\;+ \texttt{MIN\_IP\_PKT\_LEN}       & 38 & 0.997545 & $-7.3\times10^{-5}$ \\
\bottomrule
\end{tabular}
\end{table}

The near-deterministic features are therefore removable, not
load-bearing, and the natural reading that the corpus is saturated
because of a small artefact is not supported. What the ablation shows
instead is more general and, for the purposes of the study, more
limiting. With the three features withheld, the strongest remaining
single feature, \texttt{SERVER\_TCP\_FLAGS}, still separates the
classes at $0.8573$ accuracy under one threshold against a
majority-class baseline of $0.5103$, and each of the ten most important
remaining features exceeds that baseline by between $0.17$ and $0.35$.
Separability is thus distributed redundantly across many attributes
rather than concentrated in a few. A benchmark with this property
cannot discriminate between architectures on clean performance, because
any reasonable model recovers the same decision from a different subset
of features, and that is precisely the pattern the clean results show:
four of the five architectures separate by $1.8\times10^{-3}$, an
interval of the same order as the largest of their seed dispersions,
with CNN1D trailing them by a further $2.8\times10^{-3}$.

Two figures in this subsection are easily confused and are therefore
distinguished explicitly. The single-threshold accuracy of
\texttt{SERVER\_TCP\_FLAGS} is $0.8633$ when all features are present
and $0.8573$ once the three near-deterministic features are withheld;
the paper quotes the latter, because the claim it supports concerns
what remains after their removal.

%=====================================================================
\section{Padding-masked re-evaluation in full}
\label{sup:pad}

The pseudo-sequence construction pads both NSL-KDD and NF-UNSW-NB15
from $41$ to $64$ positions, so $35.94\%$ of each input vector is a
constant zero. An $\ell_{\infty}$-bounded adversary perturbs every
coordinate, including the padding, and therefore spends roughly a third
of its budget on dimensions that are constant in training. Padding
coordinates have no counterpart in any emittable flow, so evasion
attributable to them is an artefact of the representation rather than a
security finding, which makes the size of their contribution a question
about the validity of the comparison and not merely about its
interpretation.

The white-box grid was therefore repeated on the two padded corpora
with the perturbation confined to the informative coordinates, the
padding positions being held at their clean value of zero at every
step. The replication is reduced to the seeds $\{13,42\}$ for cost and
was executed on an NVIDIA A100-SXM4-40\,GB device
(\texttt{results/e3\_padding\_mask.json}). Before the masked arm is
read, the harness itself was validated against the published unmasked
figures on all $24$ (corpus, model, quantity) checks: it reproduces the
clean and FGSM values to within $8.9\times10^{-5}$ and the PGD values,
which begin from a random point in the ball, to within
$3.4\times10^{-3}$, so the two arms are comparable.

\begin{table}[h]
\centering
\caption{Padding-masked re-evaluation. Mean robust accuracy over the
white-box grid with the perturbation applied to all coordinates
(unmasked) and to the informative coordinates only (masked), averaged
over the two corpora, both attacks and four budgets.}
\label{tab:padagg}
\small
\begin{tabular}{@{}lccc@{}}
\toprule
Model & Unmasked & Masked & Change \\
\midrule
Transformer & 0.8605 & 0.8648 & $+0.0044$ \\
Mamba       & 0.6633 & 0.7286 & $+0.0652$ \\
CNN1D       & 0.6608 & 0.7500 & $+0.0892$ \\
BiLSTM      & 0.5563 & 0.7155 & $+0.1592$ \\
\bottomrule
\end{tabular}
\end{table}

The masked attack proves weaker for every architecture, which settles
the direction of the effect: perturbing the padding assists the
adversary rather than wasting its budget. The assistance is strongly
architecture dependent, as Table~\ref{tab:padagg} shows. The
Transformer is almost indifferent to the padding, its two arms on
NF-UNSW-NB15 differing by less than $2\times10^{-3}$ under either
attack, whereas the state-space detector derives a measurable part of
its apparent fragility from coordinates no adversary can reach.

The dependence is also unevenly distributed in a way a narrower
experiment would have misread. Under FGSM on NSL-KDD alone the
state-space delta is $+0.0070$, which taken by itself would suggest
padding independence; the dependence is concentrated in PGD, where it
reaches $+0.1005$ on NF-UNSW-NB15 and $+0.0998$ on NSL-KDD.

The consequence for the principal finding is that its direction
survives and its magnitude does not. On this subset the mean white-box
margin in favour of the Transformer falls from $+0.1971$ unmasked to
$+0.1363$ masked, a reduction of $31\%$, and the architecture ordering
is unchanged in both arms. No cell inverts: the margin stays positive
in all four (corpus, attack) combinations, narrowing most under PGD on
NF-UNSW-NB15, from $+0.2934$ to $+0.1947$, and not at all under FGSM on
NSL-KDD ($+0.0732$ to $+0.0735$). The attention detector therefore
remains the more robust architecture under a realistic adversary, but
by roughly two thirds of the margin the unmasked grid reports, and the
larger figure should not be quoted without that qualification.

%=====================================================================
\section{What the aggregate metrics conceal about severity}
\label{sup:severity}

Aggregate robust accuracy states how much detection is lost, not which
intrusions are lost with it, and the per-category decomposition shows
that the answer is architecture dependent and not stable across
corpora. On NSL-KDD the volumetric categories dominate for the
convolutional and recurrent baselines, with evasion rates of $0.773$
and $0.804$ for DoS and $0.757$ and $0.472$ for Probe, whereas for the
state-space and attention detectors the largest rates occur on U2R
($0.194$) and Probe ($0.167$) respectively while DoS evasion remains at
$0.068$ and $0.073$. Clean detection rates for R2L and U2R are low for
every architecture ($0.104$--$0.313$), so the corresponding evasion
rates are bounded from above by a small quantity.

On NF-UNSW-NB15 clean detection remains at or above $0.9840$ in $35$ of
the $36$ cells, the exception being CNN1D on Worms ($0.9600$), where one
instance of $25$ is missed, and under attack the four architectures
separate cleanly: the Transformer registers at most $0.004$ (DoS)
across all nine categories, Mamba's largest rates are $0.143$
(Generic), $0.125$ (Shellcode) and $0.036$ (DoS) with $0.000$ on Worms
and $0.001$ on Reconnaissance, BiLSTM spans $0.080$ (Worms) to $0.938$
(Analysis), and CNN1D is evadable in every category but Worms, reaching
$0.994$ (Analysis), $0.982$ (Shellcode) and $0.967$ (Backdoor).

On CSE-CIC-IDS2018 clean detection of the malicious categories is weak
in absolute terms, between $0.015$ (CNN1D) and $0.149$ (Mamba) for
Infiltration and between $0.000$ (CNN1D) and $0.006$ (Transformer) for
Bot, so evasion rates are correspondingly small, the largest being
$0.132$ (Mamba, Infiltration). One negative value is recorded,
$-0.494$ for BiLSTM on Bot, where the adversarial detection rate
($0.498$) exceeds the clean rate ($0.004$): the perturbation moves that
architecture from labelling almost every Bot flow benign to labelling
about half of them malicious. A reversal of this kind is what one
expects when the clean decision on a category is essentially
uninformative, and Bot is precisely the family the chronological split
places wholly outside the training interval.

A caveat concerns what the evasion success rate measures. Being defined
on malicious instances alone, it is blind to how an architecture treats
benign traffic, and the two can move in opposite directions. On
NF-UNSW-NB15 under FGSM at $\varepsilon=0.1$ Mamba retains $0.955$ of
its malicious detection on the first seed while its aggregate accuracy
on that seed falls to $0.473$, which places its accuracy on the benign
class at $0.011$: the perturbation drives it to label almost every
benign flow malicious. CNN1D fails in the opposite direction on the
same corpus, holding $0.974$ on the benign class while its malicious
detection falls to $0.140$. A low evasion rate is therefore not by
itself evidence of robustness, and the two-class decomposition released
with the artefact should be read alongside the category table.

%=====================================================================
\section{Adversarial training: effective, but insufficiently verified}
\label{sup:defence}

FGSM-mixed adversarial training raises mean white-box robust accuracy
in all six (dataset, model) combinations, by between $+0.1734$ and
$+0.2506$ for the state-space detector and between $+0.2005$ and
$+0.2498$ for the convolutional baseline, at a clean-accuracy cost that
never exceeds $0.0167$ in magnitude and is of mixed sign
(\texttt{results/defense\_results.csv}). The gains are largest at the
largest budget, where the undefended models were weakest, and for both
architectures they fall in the same cell: NF-UNSW-NB15 under PGD at
$\varepsilon=0.1$, where the state-space detector rises from $0.2508$
to $0.9762$ and the convolutional baseline from $0.3888$ to $0.9129$.

One qualification governs how far these gains can be taken, and it is
the reason the paper treats the defence arm as incompletely validated.
The defended checkpoints were evaluated under the white-box grid only;
no black-box transfer attack was run against them. Single-step
adversarial training is known to be prone to catastrophic overfitting,
in which robustness to the single-step attack rises sharply while
robustness to multi-step attacks collapses, and the resulting model
exhibits a degenerate loss surface on which a single gradient poorly
reflects the loss. Two observations argue against that pathology here:
the gains are realised under PGD as well as FGSM, and they are realised
at every budget rather than only at the one the defence was trained at.
But the decisive check is the transfer-attack check, precisely because
it does not rely on the gradient the defence may have obscured, and it
was not performed. The surrogates it requires are already trained and
released, so the experiment is inexpensive; until it is run, the
defence is validated only by the class of attack whose gradients its
training may have shaped.

%=====================================================================
\section{Threat-model realism}
\label{sup:threat}

Perturbations in this study are generated in standardised feature space
and are constrained only by the empirical feature box of the paper's
Equation~(1). They are not required to correspond to flows that an
adversary could actually emit, and the gap between the two is not
incidental. A NetFlow record's fields are linked by protocol semantics
and by arithmetic identities: byte and packet counters are
non-negative integers, duration and throughput are related by
definition, flag fields take values from a discrete set, and a
perturbation that raises the packet count without raising the byte
count describes a flow that cannot exist. An unconstrained
$\ell_{\infty}$ perturbation violates such relations freely.

Three consequences follow, and the paper states them rather than
leaving them to be inferred. The reported robust accuracies are
pessimistic as estimates of operational risk, since part of the
measured degradation is caused by inputs no adversary can produce. The
defence gains may be correspondingly optimistic, since a defence
hardened against unconstrained perturbations has been hardened partly
against a fiction. And whether the architecture ranking is stable under
a constraint-aware attack is unknown: the literature is divided on
whether domain constraints confer material robustness, with one line of
work reporting that up to $82\%$ of the adversarial examples PGD
produces violate learned constraints and that projecting onto the
constrained set restores up to $34\%$ of accuracy, and another
reporting that constraint-compliant attacks still reach
misclassification rates above $95\%$ although an adversary restricted
to five randomly chosen features, about a tenth of that space, succeeds
just under $50\%$ of the time.

Establishing which reading holds for this comparison requires either a
projection onto the feasible set after each attack step or an attack
that searches within it. Neither was performed here, and the paper's
robustness figures should be read as bounding the attainable evasion
rate from above rather than as certifying an end-to-end network-level
exploit.

%=====================================================================
\section{Capacity ablation in full}
\label{sup:capacity}

All results in the paper rest on a single capacity setting shared by
every architecture, and the clean-accuracy comparison was therefore
repeated at $d_{\mathrm{model}}\in\{32,64\}$ over the two seeds
$\{13,42\}$ with all remaining hyperparameters fixed
(\texttt{results/ablation\_results.csv}).

\begin{table}[h]
\centering
\caption{Capacity ablation. Mean clean accuracy over two seeds at two
embedding widths.}
\label{tab:cap}
\small
\begin{tabular}{@{}llccc@{}}
\toprule
Dataset & Model & $d_{\mathrm{model}}=32$ & $d_{\mathrm{model}}=64$ & Change \\
\midrule
\multirow{4}{*}{NSL-KDD}
 & Mamba       & 0.7435 & 0.7658 & $+0.0223$ \\
 & Transformer & 0.7850 & 0.7936 & $+0.0086$ \\
 & CNN1D       & 0.8005 & 0.7549 & $-0.0456$ \\
 & BiLSTM      & 0.7868 & 0.8139 & $+0.0271$ \\
\midrule
\multirow{4}{*}{NF-UNSW-NB15}
 & Mamba       & 0.9970 & 0.9971 & $+0.0001$ \\
 & Transformer & 0.9966 & 0.9967 & $+0.0001$ \\
 & CNN1D       & 0.9872 & 0.9949 & $+0.0076$ \\
 & BiLSTM      & 0.9947 & 0.9964 & $+0.0016$ \\
\midrule
\multirow{4}{*}{CSE-CIC-IDS2018}
 & Mamba       & 0.5694 & 0.5153 & $-0.0541$ \\
 & Transformer & 0.5705 & 0.5166 & $-0.0539$ \\
 & CNN1D       & 0.5040 & 0.5460 & $+0.0420$ \\
 & BiLSTM      & 0.5136 & 0.5681 & $+0.0545$ \\
\bottomrule
\end{tabular}
\end{table}

Two inversions must be recorded as qualifications of the principal
findings. The four-way ordering is capacity-invariant only on
NF-UNSW-NB15, where Mamba leads at both widths; on NSL-KDD the leading
architecture changes with width, CNN1D being strongest at $32$
($0.8005$) and BiLSTM at $64$ ($0.8139$), and on CSE-CIC-IDS2018 the
ordering inverts outright, the two highest-scoring models at $32$ being
the two lowest at $64$. With two seeds and a dispersion reaching
$\pm0.0834$ in the CSE-CIC-IDS2018 cells, that inversion is most
parsimoniously read as noise rather than as a capacity effect, which is
itself the finding: on a corpus where models do not learn, an ablation
cannot discriminate.

The pairwise relation carrying the central comparison of the study is
also not uniform. The Transformer exceeds Mamba on clean accuracy in
four of the six (corpus, width) cells and trails it in the two
NF-UNSW-NB15 cells by $0.0003$ and $0.0004$. That margin is substantial
only on NSL-KDD, at $0.0415$ at $d_{\mathrm{model}}=32$ and $0.0277$ at
$64$, in both cases larger than either model's seed dispersion; across
the four cells on the other two corpora its magnitude never exceeds
$0.0013$, well inside the dispersion observed on CSE-CIC-IDS2018, so
those four cells establish no ordering in either direction.

Two restrictions bound what the ablation can settle in any case. It
varies clean accuracy alone, and therefore offers no evidence on
whether the robustness findings, which are the actual subject of the
study, are artefacts of the chosen capacity; and with two widths and
two seeds it lacks the power to resolve the CSE-CIC-IDS2018 inversion.

%=====================================================================
\section{Index of the deposited tables and figures}
\label{sup:index}

The tables and figures the paper omits are deposited with the artefact
rather than reproduced here, in the forms below. Aggregated tables are
in \texttt{results/tables/} as CSV, Markdown and \LaTeX{}; the rendered
copies used during preparation are in \texttt{published/tables\_tex/}
and \texttt{published/figures/}, with PNG versions in
\texttt{published/figures\_png/}.

\begin{table}[h]
\centering
\caption{Deposited tables.}
\small
\begin{tabular}{@{}l p{0.62\textwidth}@{}}
\toprule
File & Contents \\
\midrule
\texttt{T1\_ana\_ozet\_sonuclari}      & All five detection metrics per (dataset, model, attack, $\varepsilon$) cell, mean and standard deviation \\
\texttt{T2\_istatistiksel\_...} & Primary paired comparison against the pre-specified comparator \\
\texttt{T2b\_..\_posthoc}              & Secondary comparison against the per-cell strongest competitor \\
\texttt{T3\_saldiri\_turu\_...} & Evasion success rate per attack category \\
\texttt{T4\_cikarim\_verimliligi}      & Parameter counts, latency and throughput at both batch regimes \\
\texttt{T5\_savunma\_etkinligi} & Defended against undefended, per (dataset, model, attack, $\varepsilon$) \\
\texttt{T6\_hiperparametre\_...} & Capacity ablation \\
\bottomrule
\end{tabular}
\end{table}

One further frame is deposited under
\texttt{results/superseded/}: the result file of the first execution of
the grid, before the \texttt{Attack} leakage of
Section~\ref{sup:leak} was corrected. No figure in the paper comes from
it. It is kept so that one claim which would otherwise rest on
testimony can be checked, namely that the obfuscated-gradient
diagnostic (5) fired before the correction and does not after it: in
that frame the Transformer's accuracy on NF-UNSW-NB15 never falls below
$0.9976$ at any budget, so a tenfold increase in $\varepsilon$ costs it
at most $0.0024$, whereas in the corrected run the same tenfold
increase costs $0.051$ under FGSM and $0.046$ under PGD. A leaking
feature can therefore imitate gradient masking on this diagnostic, and
the two are separated only by auditing the features rather than the
curves.

\begin{table}[h]
\centering
\caption{Deposited figures, in \texttt{published/figures/}.}
\small
\begin{tabular}{@{}l p{0.60\textwidth}@{}}
\toprule
File pattern & Contents \\
\midrule
\texttt{clean\_acc\_*}   & Clean accuracy by architecture, one panel per corpus \\
\texttt{robust\_fgsm\_*} & Robust accuracy against $\varepsilon$ under FGSM, one per corpus \\
\texttt{robust\_pgd\_*}  & The same under PGD \\
\texttt{catvuln\_*}      & Evasion success rate by attack category, representative cells \\
\texttt{latency\_bs1}, \texttt{latency\_bs512} & Measured latency at the two batch regimes \\
\texttt{params\_vs\_robust} & Parameter count and latency against mean robust accuracy \\
\texttt{defense\_*}      & Robustness curves before and after adversarial training \\
\texttt{ablation\_*}     & Clean accuracy against embedding width, one per corpus \\
\bottomrule
\end{tabular}
\end{table}

\noindent
The two verification records are
\texttt{results/verification/verification.json}, which holds the two numeric
S6 implementation checks, the split verification and the
single-threshold leakage audit, and
\texttt{results/verification/ttl\_ablation.json}, which holds the
feature-removal ablation of Section~\ref{sup:leak}. The padding-masked
re-evaluation of Section~\ref{sup:pad} is in
\texttt{results/e3\_padding\_mask.\{json,csv,tex\}}. The resolved run
configuration, including every seed, is in
\texttt{results/logs/config.json}.

\end{document}
